%% ma: L12-B16-0001
%% lop: 12
%% chuong: 5
%% bai: 16
%% dang: 12.16.1
%% loai: DS
%% muc: [NB, TH, VD, VD]
%% dap_an: "ĐĐSS"
%% nguon: "(NBV)-ĐỀ SỐ 10-MỖI NGÀY 1 ĐỀ THI 2025.pdf (D10, MỖI NGÀY 1 ĐỀ - Đề số 10), Phần II Câu 4"
%% kien_thuc: "Bài 15 (phương trình tham số của đường thẳng), Bài 8 (độ dài vectơ), Bài 16 (góc giữa đường thẳng và mặt phẳng); chuyển động thẳng đều"
%% trang_thai: chinh-thuc
%% ghi_chu: "Đã duyệt 11/10/2026. Đã bỏ đoạn giới thiệu công nghệ đường ray dây (đọc thêm, không chứa dữ kiện toán) và ảnh minh họa. Đề gốc ghi t thuộc R ở ý b, đã sửa thành t >= 0 cho khớp với thời gian chuyển động. Đáp án gốc Đ Đ S S đúng."
\begin{ex}
Trong không gian với hệ tọa độ $Oxyz$ (đơn vị trên mỗi trục là mét), một cabin xuất phát từ điểm $A(10;3;0)$ và chuyển động đều theo đường cáp có vectơ chỉ phương $\vec u=(2;-2;1)$ (hướng chuyển động cùng chiều với hướng vectơ $\vec u$) với tốc độ là $4{,}5$ (m/s).
\choiceTF
{\True Phương trình tham số của đường thẳng chứa đường cáp là $\begin{cases}x=10+2t\\y=3-2t\\z=t\end{cases}$ $(t\in\mathbb R)$.}
{\True Giả sử sau thời gian $t$ (s) kể từ khi xuất phát ($t\ge0$), cabin đến điểm $M$. Khi đó tọa độ điểm $M$ là $\left(3t+10;-3t+3;\dfrac{3t}{2}\right)$.}
{Cabin dừng ở điểm $B$ có hoành độ $x_B=550$, khi đó quãng đường $AB=800$ m.}
{Đường cáp $AB$ tạo với mặt phẳng $(Oxy)$ một góc $30^\circ$.}
\loigiai{a) Đường cáp đi qua $A(10;3;0)$ và có vectơ chỉ phương $\vec u=(2;-2;1)$ nên có phương trình tham số như đã cho. Đúng.

b) $|\vec u|=\sqrt{2^2+(-2)^2+1^2}=3$. Sau $t$ giây cabin đi được $4{,}5t$ m, tức $\overrightarrow{AM}=\dfrac{4{,}5t}{3}\vec u=1{,}5t\,\vec u=(3t;-3t;1{,}5t)$. Vậy $M(3t+10;-3t+3;\tfrac{3t}2)$. Đúng.

c) $x_B=550\Rightarrow3t+10=550\Rightarrow t=180$ s, nên $AB=4{,}5\cdot180=810$ m, không phải $800$ m. Sai.

d) Mặt phẳng $(Oxy)$ có vectơ pháp tuyến $\vec k=(0;0;1)$. Gọi $\varphi$ là góc giữa đường cáp và $(Oxy)$ thì $\sin\varphi=\dfrac{|\vec u\cdot\vec k|}{|\vec u|\cdot|\vec k|}=\dfrac13$, suy ra $\varphi\approx19{,}5^\circ\ne30^\circ$. Sai.}
\end{ex}
